% ===========================================================================
%  utente/09-account/00-capitolo.tex
% ===========================================================================
\chapter{Account e abbonamento}\label{cap:account}

\section{A che serve l'account}\label{sez:account-serve2}

A quattro cose, e a nessun'altra: entrare nel marketplace e ricevere
messaggi --- tutte e due nella \voce{Community} --- avere un profilo
pubblico, sottoscrivere un abbonamento. Tutto il resto dell'applicazione
funziona senza.

\figuraSchermata{0.86}{figure/schermate/subscription.png}
  {Le pagine che richiedono un account lo dicono e offrono di farlo, invece
   di sparire dal menu.}
  {subscription}

\section{Entrare la prima volta}\label{sez:entrare}

\begin{pratica}{Durante la fase di prova}
In beta l'accesso è riservato a chi si è iscritto dal sito. Il primo
ingresso non è una registrazione: è un recupero password.

\begin{enumerate}
  \item Nella pagina di accesso, scegliere \voce{Password dimenticata}.
  \item Inserire l'indirizzo con cui ci si è iscritti alla beta.
  \item Arriva un messaggio con un codice; con quello si imposta una
        password nuova.
  \item Da lì in poi si entra normalmente.
\end{enumerate}

Un indirizzo non presente nelle iscrizioni viene rifiutato, sia
all'accesso sia alla registrazione.
\end{pratica}

\noindent
Si può anche entrare con Google: l'accesso si apre nel browser di sistema e
l'applicazione raccoglie l'esito.

\section{Il profilo}\label{sez:profilo}

\figuraSchermata{1.0}{figure/schermate/profile.png}
  {Il profilo. Nome visibile, biografia, e l'uscita dall'account.}
  {profile}

\noindent
Da qui si cambiano nome visibile e biografia --- quelli che gli altri vedono
accanto a un artefatto pubblicato --- e si manda un riscontro sul prodotto.

La scheda \voce{Security} tratta password, indirizzo e verifica in due
passaggi. Non sono conservati dall'applicazione: stanno sul sito, e
l'applicazione fa solo da tramite.

\section{I tre abbonamenti}\label{sez:abbonamenti}

Servono a una cosa sola: poter \emph{pubblicare} sul marketplace. Comprare,
scaricare, installare e usare non richiedono niente.

\begin{center}
\begin{tabular}{@{}L{22mm}L{116mm}@{}}
\toprule
\textbf{Livello} & \textbf{Che cosa dà} \\
\midrule
\rigaTenue
\voce{Free} & Tutto il sistema: costruire, provare, far girare, comprare,
  installare. È un utente pieno, non una versione ridotta. \\
\voce{Author} & In più: pubblicare sul marketplace. \\
\rigaTenue
\voce{Studio} & Lo stesso, a un prezzo diverso. \\
\voce{Desk} & Lo stesso, a un prezzo diverso. \\
\bottomrule
\end{tabular}
\end{center}

\begin{nota}{Perché oggi i tre piani fanno la stessa cosa}
È voluto, ed è temporaneo. Durante la fase di prova i tre livelli danno solo
la pubblicazione e differiscono per prezzo; che cosa aggiungere ai due
superiori --- o se tenerne uno solo --- si decide all'uscita dalla beta. I
prezzi attuali sono segnaposto.
\end{nota}

\subsection{Mensile o annuale}

Entrambi, con l'annuale scontato. Si sceglie al momento dell'acquisto e si
cambia quando si vuole: il sistema calcola il conguaglio da sé.

\subsection{Disdetta, carta, fatture}

Dal pulsante del portale di fatturazione. Si apre una pagina del fornitore
dei pagamenti, dove si disdice l'abbonamento, si cambia il metodo di
pagamento e si scaricano le fatture.

\begin{nota}{La disdetta non toglie niente di già comprato}
Un artefatto acquistato resta acquistato e resta installato. L'abbonamento
riguarda la possibilità di \emph{pubblicare}, non quella di usare.
\end{nota}

\section{Quando qualcosa non torna}\label{sez:account-problemi}

\begin{center}
\begin{tabular}{@{}L{50mm}L{88mm}@{}}
\toprule
\textbf{Sintomo} & \textbf{Che cosa fare} \\
\midrule
\rigaTenue
«Non sei abbonato» ma l'abbonamento c'è
  & Chiudere e riaprire l'applicazione: il collegamento con il sito viene
    stabilito all'avvio. \\
Il marketplace è vuoto
  & Stesso motivo. Se persiste, uscire dall'account e rientrare. \\
\rigaTenue
Ho pagato e non vedo l'artefatto
  & Aspettare qualche secondo e ricaricare. Se resta così, usare il comando
    che spiega perché non risulta acquistato, nella scheda dell'artefatto
    (\vediCap{marketplace}). \\
Una pagina dice «non lo so» invece di un numero
  & Non è un guasto: l'applicazione non è riuscita a chiedere al sito, e lo
    dice invece di inventare una risposta. Riprovare più tardi. \\
\bottomrule
\end{tabular}
\end{center}
